\documentclass[10pt,a4paper]{amsart}
\usepackage[margin=19mm]{geometry}
\usepackage{amsmath,amssymb,mathtools}
\usepackage[scr=rsfs,cal=euler]{mathalfa}
\usepackage{xfrac}
\usepackage[unicode,hidelinks,bookmarks=false]{hyperref}
\newcommand{\ie}{i.e.\ }
\newcommand{\cf}{cf.\ }
\newcommand{\resp}{respectively\ }
\newcommand{\Subsection}{\subsection}
\providecommand{\nobreakdash}{\nobreak\hbox{-}}
\setlength{\emergencystretch}{2em}
\multlinegap0pt
\usepackage{siunitx}
\usepackage{siunitx}
\usepackage{siunitx}
\usepackage{siunitx}
\usepackage{amssymb}
\usepackage{tikz}
\usepackage{enumerate}
\usepackage{mathtools}
\usepackage{siunitx}
\usepackage{xcolor}
\usepackage{algorithm}\usepackage{algpseudocode}
\usepackage{tcolorbox}
\newtheorem{lemma}{Lemma}
\newtheorem{theorem}{Theorem}
\newcommand{\R}{\mathbb R}
\newcommand{\bun}{{\mathcal B}}
\newcommand{\bunr}{{\mathcal B}_{\R}}
\newcommand{\cbunr}{{\overline{\mathcal B}}_{\R}}
\newcommand{\cod}{{\rm codim}_\R}
\newcommand{\ii}{\mathfrak i}
\newcommand{\orbr}{{\mathcal O}_{\R}}
\title{Generic real Jordan canonical forms}
\author{Fernando De Ter\'an, Froil\'an M. Dopico}
\date{Source: arXiv:2601.15033v1 (CC BY 4.0)}
\begin{document}
\maketitle

\noindent\emph{Source and licence.} Generic real Jordan canonical forms. Version arXiv:2601.15033v1 (CC BY 4.0), distributed by arXiv under the Creative Commons Attribution 4.0 International licence, http://creativecommons.org/licenses/by/4.0/, as recorded in the arXiv licence metadata for this exact version and shown on the abstract page https://arxiv.org/abs/2601.15033v1. Full licence notice: \texttt{licenses/arxiv-2601.15033-CC-BY-4.0.txt}.\par\smallskip

\noindent\textit{Editorial scope.} Counted unit: the article's theorem main.th (source0.tex lines 216--229). The article's complete original proof of it is delivered with it (source0.tex lines 230--266, one \texttt{proof} environment of 37 lines). No mathematics is written, repaired or re-derived: the statement and the proof are the article's own lines, in the article's own notation, and no retained line is substituted or re-flowed.  Retained uncounted as the source-local context the proof needs: lines 164--166; lines 170--172; lines 180--182; lines 192--198.  Nothing was omitted from any retained span, so the package carries no omission record.  The retained lines cite 2 works, whose entries are transcribed verbatim from the article's own bibliography source in the e-print and delivered with short editorial labels recorded in the item's reference mapping.  No retained line contains a figure environment, an image-inclusion command or drawing code, so no image is delivered and the package declares no asset dependency.  Editorial formatting only, in addition: an AMS \texttt{amsart} preamble that restates the article's own declarations and macros used by the retained lines, namely the environments \texttt{lemma}, \texttt{theorem} on separate counters, together with the standard packages those lines need.  The retained environments print in this document as Lemma 1, Theorem 1 in order of appearance, whereas the article prints the same items with the numbering of its own full document.  The complete native source of the article as distributed in the arXiv e-print for this exact version is bundled unchanged as \texttt{sources/193130/source0.tex}.
\begin{equation} \label{eq.codimorb}
\cod\orbr(A)=\dim_\R\{X\in\R^{n\times n} \, : \, XA-AX=0  \},
\end{equation}

\begin{equation}\label{rjcf}
    {\rm RJCF}(A)=\bigoplus_{i=1}^r\left(\bigoplus_{j=1}^{d_i}C_{\ell_{j,i}}(\tilde a_i, \tilde b_i)\right)\oplus\bigoplus_{i=1}^s\left(\bigoplus_{j=1}^{e_i} J_{k_{j,i}}(\tilde c_i)\right),
\end{equation}

\begin{equation} \label{eq.codimbund}
\cod\bunr(A)=\cod\orbr(A)-(2r+s).
\end{equation}

\begin{lemma}\label{nrealevals.lemma}
 If $M\in\cbunr \left(\bigoplus_{i=1}^tC(a_i,b_i)\oplus\bigoplus_{i=1}^{n-2t}[c_i]\right)$, for some $0\leq t\leq\lfloor n/2\rfloor$, with $a_i, b_i, c_i\in\R$, $ b_i >0$, then
 \begin{itemize}
     \item[\rm (a)]$M$ has, at least, $n-2t$ real eigenvalues, counting multiplicities, and
     \item[\rm(b)] if $M$ has more than $n-2t$ real eigenvalues, counting multiplicities, then at least one of the eigenvalues of $M$ is multiple.
 \end{itemize}
\end{lemma}

\begin{theorem}\label{main.th}
    The set of real $n\times n$ matrices is equal to the following finite union of bundle closures under similarity:
    \begin{equation}\label{mainid}
    \R^{n\times n}=\bigcup_{t=0}^{\lfloor n/2\rfloor}\cbunr \left(\bigoplus_{i=1}^tC(a_i,b_i)\oplus\bigoplus_{i=1}^{n-2t}[c_i]\right),
    \end{equation}
    with $a_i,b_i,c_i\in\R$, $b_i >0$, and where $(a_i,b_i) \neq (a_{i'}, b_{i'})$ and $c_i\neq c_{i'}$ for $i\neq i'$.

    Moreover:
    \begin{itemize}
        \item[\rm(i)] $\cod\bunr \left(\bigoplus_{i=1}^tC(a_i,b_i)\oplus\bigoplus_{i=1}^{n-2t}[c_i]\right)=0$, for all $0\leq t\leq \lfloor n/2\rfloor$,
        \item[\rm(ii)]  $\bunr \left(\bigoplus_{i=1}^tC(a_i,b_i)\oplus\bigoplus_{i=1}^{n-2t}[c_i]\right)\cap\cbunr \left(\bigoplus_{i=1}^{t'}C(a_i,b_i)\oplus\bigoplus_{i=1}^{n-2t'}[c_i]\right)=\emptyset$, for $t\neq t'$, and
        \item[\rm(iii)] $\bunr \left(\bigoplus_{i=1}^tC(a_i,b_i)\oplus\bigoplus_{i=1}^{n-2t}[c_i]\right)$ is open, for all $0\leq t\leq\lfloor n/2\rfloor$.
    \end{itemize}
\end{theorem}

\begin{proof}
    Let us first prove the identity \eqref{mainid}. For this, let $A\in\R^{n\times n}$. Assume that ${\rm RJCF}(A)$ is as in \eqref{rjcf}, so there is some invertible matrix $P\in\R^{n\times n}$ such that $P^{-1}AP={\rm RJCF}(A)$. We are going to see that $A\in \cbunr \left(\bigoplus_{i=1}^{n_1}C(a_i,b_i)\oplus \bigoplus_{i=1}^{n_2}[c_i]\right)$, where $n_1=\sum_{i=1}^r\sum_{j=1}^{d_i}\ell_{j,i}$ and $n_2=\sum_{i=1}^s\sum_{j=1}^{e_i}k_{j,i}$. For this, we consider the following sequence $\{ A_m \}_{m \in \mathbb{N}}$, constructed as a perturbation of $A = P\cdot{\rm RJCF}(A)\cdot P^{-1}$:
    $$
    \begin{array}{ccl}
    A_m&=&P\left(\displaystyle\displaystyle\bigoplus_{i=1}^r\left(\bigoplus_{j=1}^{d_i}C_{\ell_{j,i}}(\tilde a_i, \tilde b_i)+\begin{bmatrix}
        C(\frac{1}{m+j},\frac{1}{m+j})\\&C(\frac{1}{2m+j},\frac{1}{2m+j})\\&&\ddots\\&&&C(\frac{1}{\ell_{j,i}m+j},\frac{1}{\ell_{j,i}m+j})
    \end{bmatrix}\right)\right.\\&&\left.\displaystyle\oplus \bigoplus_{i=1}^s\left(\bigoplus_{j=1}^{e_i} J_{k_{j,i}} (\tilde c_i) +\begin{bmatrix}
        \frac{1}{m+j}\\&\frac{1}{2m+j}\\&&\ddots\\&&&\frac{1}{k_{j,i}m+j}
    \end{bmatrix}\right)\right)P^{-1}.
    \end{array}
    $$
Note that the sequence $\{A_m\}_{m\in\mathbb N}$ converges to $A = P\cdot {\rm RJCF}(A)\cdot P^{-1}$. Moreover, $A_m\in \bunr \left(\bigoplus_{i=1}^{n_1}C(a_i,b_i)\oplus \bigoplus_{i=1}^{n_2}[c_i]\right)$ for $m$ large enough, with $c_i\neq c_{i'}$ and $(a_i , b_i) \neq (a_{i'}, b_{i'})$, for $i\neq i'$, because of the following:
\begin{itemize}
    \item The real eigenvalues of $A_m$ are $ \tilde c_i+\frac{1}{km+j}$, where $1 \leq i \leq s$, $1\leq j\leq e_i$, and $1\leq k\leq k_{j,i}$. Let us see that all of them are distinct for $m$ sufficiently large. If $\tilde c_i+\frac{1}{km+j}=  \tilde c_i+\frac{1}{k'm+j'}$, for some $1\leq k \leq k_{j,i}\,$,  $1\leq k'\leq k_{j',i}\,$, and $1\leq j,j'\leq e_i\,$, then it must be $(k-k')m=j'-j$. But, for $m$ large enough, this is not possible unless $k=k'$ and $j=j'$, because  $|j'-j|\leq e_i$. For $i\neq i'$, since $\tilde c_i\neq \tilde c_{i'}$, a value of $m$ large enough guarantees that $\tilde c_i+\frac{1}{km+j}\neq \tilde c_{i'}+\frac{1}{k'm+j'}$, for any $1\leq k\leq k_{j,i},1\leq j\leq e_i,1\leq k'\leq k_{j',i'}$, and $1\leq j'\leq e_{i'}$, because $\frac{1}{km+j}$ and $\frac{1}{k'm+j'}$ tend to $0$ as $m$ tends to infinity.
    \item Similarly, the non-real complex conjugate eigenvalues of $A_m$ are $\tilde a_i+\frac{1}{\ell m+j} \pm (\tilde b_i+\frac{1}{\ell m+j})\ii$, for $1\leq i\leq r$, $1\leq j\leq d_i$, and $1\leq \ell\leq\ell_{j,i}$. To see that all of them are distinct for $m$ large enough, it is sufficient to check that those corresponding to the ``$+$'' sign are all distinct, since $\widetilde b_i >0$. If $\tilde a_i+\frac{1}{\ell m+j}+ (\tilde b_i+\frac{1}{\ell m+j})\ii= \tilde a_{i}+\frac{1}{\ell'm+j'}+ (\tilde b_{i}+\frac{1}{\ell'm+j'})\ii$, for some $1\leq \ell \leq \ell_{j,i}\,$,  $1\leq \ell' \leq \ell_{j',i}\,$, and $1\leq j,j'\leq d_i\,$, we conclude, as before, that, for $m$ large enough, it must be $\ell=\ell'$ and $j=j'$. Also, for $i \ne i'$, since $(\tilde a_i, \tilde b_i )\neq (\tilde a_{i'}, \tilde b_{i'})$, we get that, for $m$ large enough, $\tilde a_i+\frac{1}{\ell m+j}+ (\tilde b_i+\frac{1}{\ell m+j})\ii\neq \tilde a_{i'}+\frac{1}{\ell'm+j'}+ (\tilde b_{i'}+\frac{1}{\ell'm+j'})\ii$, for all $1\leq \ell\leq\ell_{j,i},1\leq\ell'\leq \ell_{j',i'}\,,$ and $1\leq j\leq d_i,1\leq j'\leq d_{i'}$. Finally, note that $\tilde b_i+\frac{1}{\ell m+j} > 0$ for $m$ sufficiently large because $\tilde b_i > 0$.
\end{itemize}
Therefore, $A\in \cbunr \left(\bigoplus_{i=1}^{n_1}C(a_i,b_i)\oplus \bigoplus_{i=1}^{n_2}[c_i]\right)$, with $n_1$ and $n_2$ as above and $a_i,b_i,c_i\in\R$, $b_i >0$, $(a_i,b_i) \neq (a_{i'}, b_{i'})$, and $c_i\neq c_{i'}$, for $i\neq i'$. Since it must be $2n_1+n_2=n$, if we set $t=n_1$, then $A\in \cbunr \left(\bigoplus_{i=1}^tC(a_i,b_i)\oplus\bigoplus_{i=1}^{n-2t}[c_i]\right)$, and this proves \eqref{mainid}.

Now let us prove claim (ii) in the statement. Assume that there is some matrix $A\in\bunr \left(\bigoplus_{i=1}^tC(a_i,b_i)\oplus\bigoplus_{i=1}^{n-2t}[c_i]\right)\cap\cbunr \left(\bigoplus_{i=1}^{t'}C(a_i,b_i)\oplus\bigoplus_{i=1}^{n-2t'}[c_i]\right)$, for some $1\leq t,t'\leq\lfloor n/2\rfloor$.

Since $A\in\bunr \left(\bigoplus_{i=1}^tC(a_i,b_i)\oplus\bigoplus_{i=1}^{n-2t}[c_i]\right)$, then $A$ has exactly $n-2t$ real eigenvalues. Moreover, since $A\in\cbunr \left(\bigoplus_{i=1}^{t'}C(a_i,b_i)\oplus\bigoplus_{i=1}^{n-2t'}[c_i]\right)$, %there is a sequence $\{S_m\}_{m\in\mathbb N}$, with $S_m\in\bunr \left(\bigoplus_{i=1}^{t'}C(a_i,b_i)\oplus\bigoplus_{i=1}^{n-2t'}[c_i]\right)$, which converges to $A$.
by part (a) of Lemma \ref{nrealevals.lemma}, $A$ has, at least, $n-2t'$ real eigenvalues, so it must be $n-2t\geq n-2t'$, namely $t'\geq t$. If $t'>t$, then part (b) in Lemma \ref{nrealevals.lemma} implies that $A$ has at least one multiple eigenvalue, which is not the case. Therefore, it must be $t'=t$.


 Let us now prove claim (iii). For simplicity, set $\bun_t:=\bunr \left(\bigoplus_{i=1}^tC(a_i,b_i)\oplus\bigoplus_{i=1}^{n-2t}[c_i]\right)$, with $a_i, b_i$, and $c_i$ as in the statement, for $0\leq t\leq\lfloor n/2\rfloor$. Let $M\in\bun_{t_0}$, for some $0\leq t_0\leq\lfloor n/2\rfloor$. We are going to see that there is some $\varepsilon>0$ such that $B(M,\varepsilon)\subseteq \bun_{t_0}$, where $B(M,\varepsilon)=\{A\in\R^{n\times n}:\ \|M-A\|_2<\varepsilon\}$ is the (open) ball of radius $\varepsilon$ centered at $M$ (and where $\|\cdot\|_2$ denotes the spectral norm, see, for instance, \cite[Example 5.6.6]{hj13}). First, by the continuity of the eigenvalues (see, for instance, \cite[Th. D.2]{hj13}), there is some $\varepsilon_0>0$ such that all matrices in $B(M,\varepsilon_0)$ have $n$ different eigenvalues, and this implies that $B(M,\varepsilon_0)\subseteq \bigcup_{t=0}^{\lfloor n/2\rfloor}\bun_t$. Now, assume, by contradiction, that, for all $\varepsilon<\varepsilon_0$, the open ball $B(M,\varepsilon)$ is not contained in $\bun_{t_0}$, which implies, since all matrices in $B(M,\varepsilon)$ have $n$ different eigenvalues, that $B(M,\varepsilon)\cap \left(\bigcup_{t\neq t_0}\bun_t\right)\neq\emptyset$. As a consequence, $M$ belongs to the closure of $\bigcup_{t\neq t_0}\bun_t$, which is equal to $\bigcup_{t\neq t_0}\overline{\bun_t}$. Hence, $M\in\bun_{t_0}\cap\overline{\bun_t}$, for some $t\neq t_0$, which is in contradiction with part (ii).


Finally, let us prove claim (i) in the statement. For this, we strongly rely on the developments in \cite[Ch. VIII]{gant}. More precisely, the dimension of the solution space of $XA-AX=0$ depends on RJCF($A$) and, for RJCF$(A)=\bigoplus_{i=1}^tC(a_i,b_i)\oplus\bigoplus_{i=1}^{n-2t}[c_i]$, with the parameters $a_i, b_i,$ and $ c_i$ as in the statement, it is equal to
$$
\sum_{i=1}^t \dim_\R\{ X\in\R^{2\times 2} \, : \, XC(a_i,b_i)-C(a_i,b_i)X=0 \}+\sum_{i=1}^{n-2t}\dim_\R\{ x\in\R \, : \, xc_i-c_ix=0 \}.
$$
Since $\dim_\R\{ x\in\R \, : \, xc_i-c_ix=0 \} =1$ and, by a straightforward calculation, $\dim_\R\{ X\in\R^{2\times 2} \, : \, XC(a_i,b_i)-C(a_i,b_i)X=0 \}=2$, we conclude from \eqref{eq.codimorb} that
$$
\cod\orbr\left(\bigoplus_{i=1}^tC(a_i,b_i)\oplus\bigoplus_{i=1}^{n-2t}[c_i]\right)=2t+ n - 2t=n,
$$
and, as a consequence of \eqref{eq.codimbund}, $\cod\bunr\left(\bigoplus_{i=1}^tC(a_i,b_i)\oplus\bigoplus_{i=1}^{n-2t}[c_i]\right)=n-n=0$.
\end{proof}

\begin{thebibliography}{99}
\bibitem[gant]{gant} F.~R. Gantmacher. \newblock {\em The Theory of Matrices}. \newblock AMS Chelsea Publishing, New York, 1959.
\bibitem[hj13]{hj13} R.~A. Horn and C.~R. Johnson. \newblock {\em {M}atrix {A}nalysis, 2nd Edition}. \newblock Cambridge University Press, 2013.
\end{thebibliography}
\end{document}
